\section{Arithmetic traces from positive ordinates}
\label{sec:prime-trace}

A second observation model gives a quantitative test for the character labels. Let $\nu_A(\gamma)=\ord_\gamma P_A$ and $S_{\mathcal A}(m)=\sum_{\chi\in\mathcal A}\chi(m)$, with repetitions counted. The next formula transfers the classical prime-power trace of Landau \cite{gonek} to the zeros of the coupled determinant. Its leading term is independent of the Hermitian matrix.

\begin{theorem}\label{thm:landau}
Under the Dirichlet critical-line input, for each fixed integer $m\ge2$ and $T\ge3$,
\[
 \sum_{0<\gamma\le T}\nu_A(\gamma)m^{i\gamma}
 =-\frac{T\Lambda(m)}{2\pi\sqrt m}\,S_{\mathcal A}(m)
  +O_{A,\mathcal A,m}\!\left(\frac{T}{\log T}\right).
\]
\end{theorem}
\begin{proof}
Use $C_A(s)=P_A((s-1/2)/i)$, $s=2+it$, $b=3/2$ and $K_b(x)=b/(b^2+x^2)$. The genus-one Hadamard product and the critical-line location of all zeros imply
\[
 \Re\frac{C_A'(2+it)}{C_A(2+it)}=c_0+\sum_\gamma\nu_A(\gamma)K_b(t-\gamma)
\]
for a real constant $c_0$. The real parts of the regularizing terms $1/(1/2+i\gamma)$ are summable. The coupled count comparison and the unit-interval zero count for each primitive completion give $\sum_{u\le\gamma\le u+1}\nu_A(\gamma)\ll_{A,\mathcal A}\log(2+|u|)$. Endpoints can be included by slightly enlarging the interval.

Put $\omega=\log m$. The Fourier integral of $K_b$ at $\omega$ is $\pi e^{-b\omega}$, by residues. Integrating the preceding identity against $e^{i\omega t}$ for $0\le t\le T$ gives
\begin{equation}\label{eq:poisson-trace}
 \int_0^T e^{i\omega t}\Re\frac{C_A'(2+it)}{C_A(2+it)}\,dt
 =\pi e^{-b\omega}\sum_{0<\gamma\le T}\nu_A(\gamma)e^{i\omega\gamma}
 +O_{A,\mathcal A,m}(\log^2(2T)).
\end{equation}
Indeed, for $0<\gamma\le T$ the removed kernel tails have mass $O((1+\gamma)^{-1}+(1+T-\gamma)^{-1})$. For exterior ordinates in $[-2T,2T]$ the retained integral is $O((1+\operatorname{dist}(\gamma,[0,T]))^{-1})$. Summing over unit intervals costs $O(\log^2(2T))$. Farther away, the retained mass is $O(T/\gamma^2)$ with total $O(\log(2T))$. These bounds also justify integrating the sum. The constant integrates to $O_m(1)$.

On the same vertical line for large $t$, logarithmically differentiate the product of completions and the determinant in \eqref{eq:model}. The determinant term is $\operatorname{tr}(M^{-1}\partial_s M)$, with $M=A+\diag(g_i)$. We have $M=(\tfrac i2\log(s/2))I+O(1)$ and $\partial_s M=O(1)$ by the differentiated Euler series and the differentiated Stirling estimate. Consequently
\[
 \frac{C'_A}{C_A}(s)=\frac{n_A}{2}\log(s/2)+c_1
 -\sum_{k\ge2}\Lambda(k)S_{\mathcal A}(k) k^{-2-it}
 +O_{A,\mathcal A}(1/\log t).
\]
The polynomial factors in $\xi$ contribute $O(1/t)$. The integral of the real part of the logarithmic main term times $e^{i\omega t}$ is $O_m(\log(2T))$ by integration by parts. In the real part of the Euler sum only the negative frequency $-\log m$ resonates. All other summands have bounded integrals after twisting, uniformly in $k$ for fixed $m$, since the nonzero values of $|\log m\pm\log k|$ have a positive lower bound. Absolute convergence yields a total error $O_{\mathcal A,m}(1)$ from these terms. The integral is therefore $-T\Lambda(m)S_{\mathcal A}(m)/(2m^2)+O_{A,\mathcal A,m}(T/\log T)$. Substitution into \eqref{eq:poisson-trace} proves the formula.
\end{proof}

\begin{corollary}[Sublinear corruption of the positive spectrum]\label{cor:sublinear}
The positive ordinates alone determine the primitive label multiset even after $o(T)$ insertions and deletions through height $T$. In particular, for any two Hermitian couplings, writing $D(T)=\sum_{0<\gamma\le T}|\nu_A(\gamma)-\nu_B(\gamma)|$, one has
\[
 \liminf_{T\to\infty}\frac{D(T)}{T}\ \ge\
 \sup_{p\text{ prime}}\frac{\log p}{2\pi\sqrt p}\,
       |S_{\mathcal A}(p)-S_{\mathcal B}(p)|.
\]
The bound is positive whenever the label multisets differ, regardless of their sizes and conductors.
\end{corollary}
\begin{proof}
Each changed multiplicity alters a trace sum by a quantity of modulus at most one. Sublinear corruption preserves every normalized trace limit. Subtract the two trace formulas at a fixed prime, take absolute values and then take the supremum. If all prime contrasts vanish, use primes in each unit residue class of a common modulus and character orthogonality, as in Theorem~\ref{thm:blind}.
\end{proof}
The comparison here counts both excesses and deficits on the positive half-line. Its leading coefficients recover the labels. The coupling coefficients are recovered by the separate inverse theorems above.
